\section{Dataset Detail}
\label{appendix:dataset_detail}

This section presents a comprehensive statistical overview of the \textbf{Deep and Broad datasets} we use in this paper, including detailed statistics metadata and licensing information, as summarized in Table~\ref{tab:dataset_stats}. Additionally, we provide individual descriptions of each dataset to elucidate their respective characteristics and intended use cases.

\begin{table*}[t]
\centering
\caption{Statistics of Deep Reasoning and Broad Reasoning Datasets. Metrics abbreviations: Comp. (Comprehensiveness), Div. (Diversity), Emp. (Empowerment).}
\scriptsize
\begin{tabular}{lcccccc}
\toprule
Dataset & Corpus Size & Chunks & Entities/Relations & Communities & Metrics & License \\
\midrule
\rowcolor{gray!20}
\multicolumn{7}{l}{\textbf{Deep Reasoning Tasks}} \\
HotpotQA & 9,809 & 9,812 & 37,358/30,987 & 5,041 & EM, F1 & Apache-2.0 \\
2WikiMultihopQA & 6,119 & 6122 & 19,311/21,077 & 3,417 & EM, F1 & Apache-2.0 \\
Musique & 11,254 & 11,300 & 32,842/39,134 & 6,258 & EM, F1 & CC-BY-4.0 \\
\midrule
\rowcolor{gray!20}
\multicolumn{7}{l}{\textbf{Broad Reasoning Tasks}} \\
CS & 10 & 2,134 & 3,530/33,507 & 1,166 & \multirow{4}{*}{Comp., Div., Emp.} & \multirow{4}{*}{Apache-2.0} \\
Agriculture & 12 & 2,025 & 6,043/12,571 & 1,039 &  &  \\
Legal & 94 & 5,900 & 26,180/44,334 & 1,359 &  &  \\
Mix & 61 & 658 & 2,784/5,089 & 425 &  &  \\
\bottomrule
\end{tabular}
\label{tab:dataset_stats}
\end{table*}

\subsection{Deep Reasoning Datasets}
\begin{itemize}
    \item \textbf{HotpotQA}~\citep{yang2018hotpotqa}: A crowdsourced question answering dataset built on English Wikipedia, comprising approximately 113K questions. Each question is constructed to require the combination of information from the introductory sections of two Wikipedia articles for answering. The dataset provides two gold paragraphs per question, along with a list of sentences identified as supporting facts necessary to answer the question. HotpotQA includes various reasoning strategies such as bridge questions (involving missing entities), intersection questions (e.g., “what satisfies both property A and property B?”), and comparison questions (comparing two entities through a common attribute). It is available in two settings: a \textit{few-shot distractor setting} where models are provided with 10 paragraphs including the gold ones, and an \textit{open-domain full-wiki setting} where models must retrieve relevant passages from the entire Wikipedia corpus given only the question.
    \item \textbf{2WikiMultihopQA}~\citep{2wiki}: A multi-hop question answering dataset that contains complex questions requiring reasoning over multiple Wikipedia paragraphs. Each question is designed to necessitate logical connections across different pieces of information to arrive at the correct answer.
    \item \textbf{Musique}~\citep{musique}: A challenging multi-hop QA dataset containing approximately 25K 2–4 hop questions, constructed by composing single-hop questions from five existing single-hop QA datasets. It is designed to feature diverse and complex reasoning paths, requiring models to integrate information from multiple hops to generate correct answers. The dataset emphasizes comprehensive evaluation of multi-step reasoning capabilities.
\end{itemize}

\subsection{Broad Reasoning Datasets}
The following datasets are curated from the UltraDomain \citep{qian2025memorag} benchmark. The benchmark construction leverages financial reports, legal contracts, and 428 college textbooks across 18 distinct domains to evaluate model versatility and adaptability in specialized and broad application scenarios:
\begin{itemize}
    \item \textbf{CS}: Computer science domain focusing on data science, software engineering, and programming topics, requiring technical comprehension and analytical reasoning.
    \item \textbf{Agriculture}: Covers agricultural practices including beekeeping, crop production, and disease prevention, demanding domain-specific knowledge integration.
    \item \textbf{Legal}: Derived from legal contracts and documents, focusing on corporate legal practices, regulatory compliance, and governance, requiring precise interpretation of nuanced legal language.
    \item \textbf{Mix}: Contains diverse contexts from college textbooks spanning natural sciences, humanities, and social sciences, testing generalization capabilities across interdisciplinary topics.
\end{itemize}